% Proposed model. Numbers in this file come from results/e1_variants.json (settings chosen on dev)
% and results/e4_efficiency.json (vocabulary sizes).
% Needs in the preamble: tikz with \usetikzlibrary{arrows.meta,positioning}, booktabs.

\section{Proposed Model}
\label{sec:model}

HinglishKey replaces every word of a tweet by the name of its \emph{equivalence class}
\citep[Section~2.2]{manning2008iir}: a term shared by all spellings that should count as the same
word. The same terms are used twice, as index terms for retrieval and as features for
classification, so one normalisation step serves both tasks (Figure~\ref{fig:pipeline}). No
lexicon, transliteration table or language identifier is used; everything except a short list of
letter rules is learned from the training tweets.

\begin{figure}[t]
\centering
\begin{tikzpicture}[
  node distance=5mm and 6mm,
  box/.style={draw, rounded corners=2pt, align=center, font=\footnotesize,
              minimum height=8mm, inner sep=3pt, text width=24mm},
  learn/.style={box, dashed},
  use/.style={box, fill=black!6},
  arrow/.style={-{Stealth[length=2mm]}, semithick}]
  \node[box] (tweet) {Raw tweet\\\texttt{Bohot acha laga!!}};
  \node[box, right=of tweet] (tok) {Tokenise\\\scriptsize lowercase, strip URLs/@, squash};
  \node[box, right=of tok] (hkey) {H-Key\\\scriptsize letter rules, vowel policy};
  \node[box, right=of hkey] (cls) {Variant class\\\scriptsize canonical spelling};
  \node[learn, below=of hkey] (fit) {Learn classes from train\\\scriptsize blocking + trigram cosine $\geq\theta$};
  \node[use, below=12mm of cls, xshift=-14mm] (index) {Positional inverted index\\\scriptsize AND, phrase, lnc.ltc};
  \node[use, right=4mm of index] (clf) {Classifiers\\\scriptsize NB, Rocchio, $k$NN};
  \draw[arrow] (tweet) -- (tok);
  \draw[arrow] (tok) -- (hkey);
  \draw[arrow] (hkey) -- (cls);
  \draw[arrow, dashed] (fit) -| (cls);
  \draw[arrow] (cls.south) -- ++(0,-4mm) -| (index.north);
  \draw[arrow] (cls.south) -- ++(0,-4mm) -| (clf.north);
\end{tikzpicture}
\caption{The HinglishKey pipeline. Each word becomes the term of its equivalence class; the
classes are learned from the training tweets only (dashed). The same terms feed the retrieval
index (E3, E4) and the classifiers (E2). The four normalisation modes stop at different points:
\texttt{raw} after tokenisation, \texttt{soundex} and \texttt{hkey} after a phonetic key, and
\texttt{hkey\_clusters} after the variant class.}
\label{fig:pipeline}
\end{figure}

\paragraph{Tokenisation.} Tweets are lowercased and split into words. URLs and @mentions are
removed, including the space-separated forms in which the SentiMix files store them, hashtags keep
their text, and runs of three or more identical letters are shortened to two
(\texttt{achhaaaa}$\to$\texttt{achhaa}).

\paragraph{H-Key.} Soundex \citep[Section~3.4]{manning2008iir} maps letters to six digit classes
designed for English surnames, and the lecture itself calls it ``not very useful'' for retrieval,
pointing to \citet{zobel1996phonetic}. Romanised Hindi varies along different lines: aspiration is
written or not (\texttt{thik}/\texttt{theek}), letters with the same sound are swapped
(\texttt{wala}/\texttt{vala}, \texttt{zindagi}/\texttt{jindagi}), consonants are doubled or not
(\texttt{accha}/\texttt{acha}), and vowels are written freely (\texttt{bahut}/\texttt{bohot}/\texttt{bht}).
H-Key is a Soundex-shaped key with one rule per kind of variation, which keeps letters instead of
digits and has no length limit by default. Two vowel policies are possible: \emph{V-drop} deletes
every vowel after the first letter, as Soundex does; \emph{V-norm} only unifies long and short
spellings of the same vowel (\texttt{ee}$\to$\texttt{i}, \texttt{oo}$\to$\texttt{u}, \ldots).
Section~\ref{sec:methodology} gives every rule.

\paragraph{Variant classes.} A phonetic key alone either merges too much (V-drop also joins
\texttt{bas}/\texttt{bus}) or too little (V-norm keeps \texttt{bahut} and \texttt{bohot} apart).
HinglishKey therefore learns classes from the training corpus: two words are merged only if they
share an H-Key \emph{and} their character-trigram vectors are similar,
\begin{equation}
\mathrm{key}(a)=\mathrm{key}(b)\ \wedge\ \cos(\vec a_{3g},\vec b_{3g})\ge\theta,
\label{eq:merge}
\end{equation}
and each class is named after its most frequent member. The key acts as a blocking step, so only
words inside one key block are ever compared.

\paragraph{Modes.} Table~\ref{tab:modes} lists the normalisation modes compared in the experiments.
All settings were chosen on the E1 development groups before any other experiment was run.

\begin{table}[t]
\centering
\caption{Normalisation modes. Vocabulary: distinct terms over all 20{,}000 SentiMix tweets (E4).}
\label{tab:modes}
\small
\begin{tabular}{@{}llr@{}}
\toprule
Mode & Term of a word & Vocabulary \\
\midrule
\texttt{raw} & the token itself & 37{,}485 \\
\texttt{soundex} & lecture Soundex code (baseline) & 4{,}756 \\
\texttt{hkey} & H-Key, V-norm (main) & 30{,}863 \\
\texttt{hkey\_clusters} & variant class: V-drop key, $\theta=0.85$ & 30{,}852 \\
\texttt{hkey\_vdrop} & H-Key, V-drop (ablation) & 15{,}709 \\
\bottomrule
\end{tabular}
\end{table}

\paragraph{Uses.} For retrieval, the terms are indexed in a positional inverted index that supports
Boolean AND, phrase queries and ranked lnc.ltc retrieval; a query is normalised with the same mode,
so a query written \texttt{bahut accha} can find a tweet written \texttt{bohot acha}. For
classification, the terms replace the words as features of multinomial Naive Bayes, Rocchio and
$k$NN. A web demo (Search, Phrase search, Variant explorer, Sentiment, Results) runs on precomputed
indexes on a laptop CPU.
